For real states $a,b$ the flux $a^{\mathsf t}C^\mu b$ has divergence $(\mathcal P a)^{\mathsf t}b+a^{\mathsf t}\mathcal P b$ by the product rule and symmetry, the bilinear conservation identity of \cite{r-25}. Its value at $(w_\ell,\delta w_\ell)$ is the derivative of $q^\mu$ in the $\delta w_\ell$ direction. In the tube coordinates $\mathsf X-\mathrm i q$ (compare \cite{r-6}, Section~1), direct differentiation shows that $D_{\ell a}$ annihilates each component of $\mathsf X-\mathrm i q$. Hence $D_{\ell a}\overline G=0$, and conjugation gives $\partial_{w_{\ell a}}G=\mathrm i V_{\ell a}G$. Each $D_{\ell a}$ has zero divergence since its state coefficients are constant and its spatial coefficients have no $\mathsf X$ dependence. Thus
\[
\operatorname{div}\mathcal Z_\ell=2\mathrm i\sum_aV_{\ell a}G_a
=2\mathrm i w_\ell^{\mathsf t}\mathcal P G=0.
\]
Here $\mathcal P G=0$ follows by differentiation in $\mathsf X$ and conjugation, since the matrices are real.

For the ordinary flat connection, $D_{\ell a}\overline G=0$ leaves only $D_{\ell a}$ acting on the coefficients of $\overline{D_{mb}}$ in the following computation:
\[
\nabla_{\mathcal Z_\ell}\overline{\mathcal Z_m}
=-\mathrm i\delta_{\ell m}\sum_\mu(G^{\mathsf t}C^\mu\overline G)\partial_{\mathsf X^\mu}.
\]
Indeed the derivative of the coefficient $(C^\mu w_m)_b$ in direction $w_{\ell a}$ is $\delta_{\ell m}(C^\mu)_{ba}$. The vector multiplying $-\mathrm i\delta_{\ell m}$ is real by the symmetry of $C^\mu$. The product rule gives $\operatorname{div}(\mathcal Z_\ell\otimes\overline{\mathcal Z_\ell})=(\operatorname{div}\mathcal Z_\ell)\overline{\mathcal Z_\ell}+\nabla_{\mathcal Z_\ell}\overline{\mathcal Z_\ell}$, whose real part vanishes. Writing $\mathcal Z_\ell=A_\ell+\mathrm iB_\ell$ with real fields exhibits the summand as $A_\ell\otimes A_\ell+B_\ell\otimes B_\ell$ after taking real parts. The imaginary part of the same flat-connection identity gives $[A_\ell,B_\ell]=\sum_\mu(G^{\mathsf t}C^\mu\overline G)\partial_{\mathsf X^\mu}$.
